% Draft v2 of the related-work paragraph (expanded to position the contribution),
% removed from main_body_compact.tex on 2026-11-16 for rework. Bib entries it uses are
% already in iclr2026/accentuation_rmt_iclr2026.bib.

\paragraph{Related work.}
{\color{red}\textit{Neural prediction and its theory.}}
Encoding models built on deep networks now predict cortical responses so similarly that benchmarks no longer rank them \citep{yamins2014performanceOptimized,guclu2015deepGradientComplexity,cadena2019deepConvolutionalV1,schrimpf2018brainscore,conwell2024inductiveBiases}. Read one way this is convergence \citep{huh2024platonicRepresentationHypothesis,kapoor2026convergentLearning}; read another it is a weak test, answered by constraining the mapping \citep{prince2024sparsePositiveAlignment} or by seeking stimuli that split models apart \citep{golan2020controversialStimuli,feather2023modelMetamers,rajalingham2018large,geirhos2019imagenetTextureBias}. {\color{red}Kernel and ridge regression theory explains which representations predict neural data through spectral bias and task--model alignment \citep{canatar2021spectralBias,canatar2023spectralTheory}; like the benchmarks, it evaluates a model on held-out samples from its training distribution.}
{\color{red}\textit{Stimulus optimization and control.}}
Closed-loop control is one limit of that strategy: the model writes the stimulus and the neuron grades it \citep{ponce2019evolvingImages,xiao2020xdream,bashivan2019neuralPopulationControl,walker2019inceptionLoops,prince2026parametric}. {\color{red}It builds on feature visualization \citep{olah2017feature,hamblin2024feature,fel2023magnitudeConstrainedOptimization}. These experiments have produced striking phenomena but little theory: backbones and synthesis procedures are chosen empirically, and why a good predictor can be a poor controller remains open \citepPW. The same hazard, that optimizing an input against a learned model exploits its errors away from the data, appears in adversarial attack \citep{nguyen2015deepNeuralFooled,goodfellow2015explainingAdversarial,madry2018towardsResistant}, where robustness has been tied to perceptually aligned gradients and to brain alignment \citep{santurkar2019imageSynthesisRobust,guo2022adversariallyTrainedNeuralRepresentations,dapello2020simulatingV1,taori2020measuring}, and in offline model-based optimization \citep{trabucco2021conservative}; there it is handled empirically or algorithmically rather than computed.}
{\color{red}\textit{Distribution shift.}}
Classic theory characterizes ridge risk under anisotropic covariance and signal--spectrum alignment \citep{dobriban2018high,hastie2022surprises,mei2022generalization,mel2021highDimensionalRegression}, as well as generalization and optimal regularization under an externally specified train--test distribution shift \citep{canatar2021outOfDistribution,tripuraneni2021covariateShift,patil2024optimalOODRidge}. {\color{red}Performative prediction studies shifts in which the data distribution responds to a deployed predictor \citep{perdomo2020performative}, through a population-level response map rather than along the predictor's own input gradient.
\textit{Our position.} We study a self-generated shift: the test distribution is supported on an optimization path created by the fitted predictor itself. Carrying the deterministic-equivalent theory of ridge regression onto this path connects the regression theory of neural prediction with feature visualization and neural control, and makes explicit when held-out prediction does and does not constrain control.}
